% =====================================================================
% Clase -- Geometría 3° -- Trimestre I -- Semana 04
% Sesión 004 (vie 25 sep., 11:00, 50 min)
% Tema: Líneas · Subtema: Líneas abiertas y cerradas; dentro, fuera y borde
% Hilo conductor del trimestre I: «El renacuajo paseador» (Rafael Pombo)
% Tema 1 de la guía: Líneas (tema:lineas)
% Material del docente (no se entrega a los estudiantes).
% =====================================================================
\documentclass[11pt]{article}
\makeatletter\def\input@path{{../../../../../plantillas/estilo/}}\makeatother
\usepackage{mmcantillo}
\guiadelestudiante{../../guia-didactica/guia-periodo-I-del-punto-al-cuerpo}

\tipodocumento{Clase}
\asignatura{Geometría}
\titulo{Semana 4: el charco de Rinrín}
\grado{3°}
\periodo{I}

% DBA: matematicas grado 3 · DBA 6 — Describe y representa formas bidimensionales y
% tridimensionales de acuerdo con las propiedades geométricas. Evidencia de esta sesión:
% «Clasifica y representa formas bidimensionales tomando en cuenta sus características
% geométricas comunes y describe el criterio utilizado» (abiertas y cerradas), y el paso a
% que una línea cerrada separa dentro, fuera y borde.

\begin{document}

\encabezado*

\begin{center}
{\renewcommand{\arraystretch}{1.3}
\begin{tabularx}{\linewidth}{@{}l l l X l@{}}
  \toprule
  \textbf{Sesión} & \textbf{Fecha} & \textbf{Min} & \textbf{Subtema} & \textbf{Entregables}\\
  \midrule
  004 & vie 25 sep., 11:00 & 50 & Líneas abiertas y cerradas; dentro, fuera y borde & Taller\\
  \bottomrule
\end{tabularx}}
\end{center}

Segundo paso del hilo. La semana pasada Rinrín dejó caminos rectos y curvos en la arena. Hoy
importa otra cosa: si el camino \emph{vuelve} a donde empezó. Y cuando vuelve, pasa algo
mágico — el papel queda partido en tres. Referencia: \refguia{tema:lineas}.

% ---------------------------------------------------------------------
\seccion{Sesión 004 — viernes 25 de septiembre · 11:00 · 50 min}

\textbf{Objetivo.} Distinguir líneas abiertas de cerradas con un criterio propio, y reconocer
que una línea cerrada separa el dentro, el fuera y el borde.

\textbf{Materiales.} Tablero; una cuerda o lana larga; hoja y colores. Ojalá un lazo o un aro.

\begin{center}
{\renewcommand{\arraystretch}{1.3}
\begin{tabularx}{\linewidth}{@{}l l X@{}}
  \toprule
  \textbf{Momento} & \textbf{Min} & \textbf{Actividad}\\
  \midrule
  Con la cuerda & 0--10 & Estirar la cuerda en el piso: abierta. Juntar las dos puntas:
    cerrada. Preguntar quién puede quedar «adentro». Que un estudiante se pare dentro del
    lazo y otro fuera. Todavía no se nombra nada: primero se ve.\\
  El nombre & 10--22 & Ahora sí: \textbf{cerrada} si terminas donde empezaste sin levantar
    el lápiz; \textbf{abierta} si no. Dibujar cuatro caminos en el tablero y que el curso
    los clasifique diciendo el porqué («empieza aquí y termina allá»).\\
  Las tres partes & 22--36 & Con el charco de Rinrín (una línea cerrada): hay cosas
    \emph{dentro}, cosas \emph{fuera} y cosas justo \emph{en el borde}. Poner cuatro cosas
    en el dibujo y preguntar una por una. La del borde es la que genera discusión: no está
    ni adentro ni afuera.\\
  Taller & 36--47 & Taller individual (5 puntos), con dibujos.\\
  Cierre & 47--50 & Recoger. Anuncio: la próxima clase, las figuras que hacen esas líneas
    cerradas cuando son rectas — y cuántos lados tiene cada una.\\
  \bottomrule
\end{tabularx}}
\end{center}

\begin{nota}
\textbf{La confusión típica.} Muchos creen que un camino es cerrado porque «se acabó la hoja»
o porque las puntas quedaron cerca. El criterio es otro y conviene repetirlo con las mismas
palabras toda la clase: \textbf{¿terminaste donde empezaste?} Es la pregunta 3 del taller.
\end{nota}

\begin{nota}
\textbf{Lo del borde.} A esta edad cuesta aceptar que algo «no esté ni adentro ni afuera».
Ayuda el ejemplo del andén o de la línea de la cancha: el pie que pisa la raya no está ni en la
calle ni en la acera. Vale la pena dejarlos discutirlo un minuto antes de resolverlo.
\end{nota}

\end{document}
